\documentclass[11pt,a4paper]{article}
\usepackage[T1]{fontenc}
\usepackage{amsmath,amssymb}
\usepackage[margin=1in]{geometry}

\title{Passthrough-Style Simulation Equations}
\author{Quest Shooting Demo}
\date{}

\newcommand{\clamp}{\operatorname{clamp}}
\newcommand{\smoothstep}{\operatorname{smoothstep}}
\newcommand{\fract}{\operatorname{fract}}
\newcommand{\lerp}{\operatorname{lerp}}

\begin{document}
\maketitle

\section{Overview}
This document describes the project's screen-space shader and its head-motion
controller. It modifies a rendered
virtual scene; it does not access the Quest camera feed. The effect is an
artistic approximation, not a calibrated model of Quest hardware.

All vector operations below are component-wise unless stated otherwise.
Screen coordinates $\mathbf{u}=(u,v)$ range from $0$ to $1$, and $t$ is
elapsed time in seconds.

\section{View center and edge weight}
Let $\mathbf{c}$ be the projected forward view center for the current rendering
view, obtained from the projection matrix in the vertex shader. For each pixel,
\begin{align}
  \mathbf{p} &= \mathbf{u}-\mathbf{c}, &
  r &= \lVert \mathbf{p}\rVert_2, &
  e &= \smoothstep(0.2,0.75,r).
\end{align}
The edge weight $e$ increases from zero toward one away from the view center.

\section{Geometric warp}
With warp strength $k_w=0.035$ and edge warp boost $k_e=1.4$ by default,
\begin{align}
  w &= 1+k_w r^2(1+k_e e),\\
  \mathbf{u}_0 &= \mathbf{c}+\mathbf{p}w,\\
  \Delta\mathbf{u}_{\text{edge}} &= 0.6\,k_w e\mathbf{p}.
\end{align}
The edge displacement is added after the swim offsets.

\section{Head motion and animated swim}
The controller compares camera orientations on consecutive frames:
\begin{align}
  R_{\Delta} &= R_{\text{previous}}^{-1}R_{\text{current}},\\
  a &= \lVert\operatorname{Euler}(R_{\Delta})\rVert_2,\\
  m &= \clamp(a\,s,0,1),\\
  h_{\text{new}} &= \lerp\!\left(
    h_{\text{previous}},m,1-e^{-\Delta t/\tau}\right).
\end{align}
Here $s=4.0$ is the default motion sensitivity, $\tau=0.12$ seconds is
the default smoothing time, and $\Delta t$ is the frame duration. The
angle change is measured per frame, rather than divided by $\Delta t$;
consequently the response also depends on frame rate.

With swim strength $k_s=0.018$ by default, $h=h_{\text{new}}$, and
$T=1.7t$, the shader calculates:
\begin{align}
  A &= k_s(0.4+1.6h),\\
  u_1 &= u_{0x}+A(0.3+e)\sin(12u_{0y}+T),\\
  v_1 &= u_{0y}+A(0.3+0.8e)\cos(11u_1+0.9T),\\
  \mathbf{u}' &= (u_1,v_1)+\Delta\mathbf{u}_{\text{edge}}.
\end{align}
The vertical offset uses the already modified horizontal coordinate $u_1$.
The $0.4k_s$ term causes swim even when the head is still. Motion here
represents rotation only, not translation or scene depth.

\section{Softness}
Let $\mathbf{D}=(W,H)$ be the viewport size in pixels, $b=0.7$ the default
softness in pixels, and $\mathbf{d}=(b/W,b/H)$. Let $I(\mathbf{x})$ be the
screen texture sampled at $\mathbf{x}$ after clamping each coordinate to
$[0.001,0.999]$. The blurred color is
\begin{align}
  B={}&0.4I(\mathbf{u}')\notag\\
  &+0.15I(\mathbf{u}'+(d_x,0))
   +0.15I(\mathbf{u}'-(d_x,0))\notag\\
  &+0.15I(\mathbf{u}'+(0,d_y))
   +0.15I(\mathbf{u}'-(0,d_y)).
\end{align}
The five weights sum to one.

\section{Grain}
The luminance of the blurred color $B=(B_R,B_G,B_B)$ is
\begin{equation}
  L=0.2126B_R+0.7152B_G+0.0722B_B.
\end{equation}
The shader constructs a pseudo-random noise value using
\begin{align}
  \mathbf{q} &= \lfloor\mathbf{u}\mathbf{D}\rfloor+\lfloor24t\rfloor,\\
  N &= \fract\!\left(
    \sin\!\bigl(\mathbf{q}\cdot(12.9898,78.233)\bigr)
    \,43758.5453\right)-0.5.
\end{align}
The scalar $\lfloor24t\rfloor$ is added to both components of $\mathbf{q}$.
With grain strength $k_g=0.015$ by default, the resulting color is
\begin{equation}
  G=B+N\,k_g\bigl(1.25-\clamp(L,0,1)\bigr).
\end{equation}
The noise is stronger in darker rendered regions; it does not measure
physical lighting conditions.

\section{Edge darkening and output}
With edge-darkening strength $k_v=0.2$ by default,
\begin{align}
  V &= \smoothstep(0.45,0.71,r),\\
  C_{\text{out}} &= \clamp\bigl(G(1-k_vV),0,1\bigr).
\end{align}
The final clamp is applied separately to each RGB component. The vignette
suggests a limited view but does not change the headset's physical field
of view.

\section{Pipeline}
\begin{equation*}
  \boxed{\text{Output}=
  \text{edge darkening}\bigl(
  \text{softened, warped scene}+\text{grain}\bigr)}
\end{equation*}
The screen texture contains opaque scene rendering; transparent objects and
UI may not receive the effect.

\end{document}
